\documentclass[12pt,a4paper]{article}
\usepackage[utf8]{inputenc}
\usepackage[portuguese]{babel}
\usepackage[T1]{fontenc}
\usepackage{amsmath}
\usepackage{amsfonts}
\usepackage{amssymb}
\usepackage{graphicx}
\usepackage{hyperref}
\usepackage{geometry}

\geometry{left=3cm,right=2cm,top=3cm,bottom=2cm}

\title{Relatório de Projeto: Rei Julian \\ \large Sistema de Rastreamento em Tempo Real}
\author{Hiago Dona Muniz}
\date{Abril de 2026}

\begin{document}

\maketitle

\noindent \textbf{Disciplina:} TEC IV \\
\textbf{Aluno:} Hiago Dona Muniz

\hrule
\vspace{0.5cm}

\section{Introdução}
O projeto "Rei Julian" consiste em uma aplicação simplificada para rastreamento geográfico em tempo real. O sistema simula o movimento de um objeto (ou pessoa) partindo de um ponto inicial fixo (Praça Coronel Pedro Osório, Pelotas-RS) e atualiza sua posição via interface web utilizando um mapa interativo.

\section{Descrição das Operações}

\begin{enumerate}
    \item \textbf{Simulação de Movimento:} O backend executa uma rotina em segundo plano que calcula uma nova coordenada a cada intervalo de 4 a 6 segundos. O cálculo utiliza trigonometria esférica para deslocar o ponto em 10 metros em uma direção aleatória, garantindo uma simulação contínua de deslocamento.
    
    \item \textbf{Comunicação em Tempo Real (SSE):} Para a transmissão de dados, foi adotada a tecnologia \textit{Server-Sent Events} (SSE). Esta escolha permite que o servidor mantenha um canal de comunicação unidirecional aberto com o cliente, enviando a nova posição em formato JSON assim que ela é gerada, sem a necessidade de requisições repetitivas por parte do navegador (\textit{polling}).
    
    \item \textbf{Interface de Visualização:} O frontend utiliza a biblioteca \textbf{Leaflet} integrada ao \textbf{OpenStreetMap} para a renderização cartográfica. Ao receber a nova coordenada via SSE, o navegador atualiza a posição do marcador no mapa e centraliza a visualização automaticamente, proporcionando uma experiência de rastreamento fluida.
\end{enumerate}

\section{Limites Potenciais e Melhorias}

\begin{itemize}
    \item \textbf{Persistência de Dados:} Atualmente, o sistema armazena a posição apenas em memória volátil. Caso o servidor seja reiniciado, o histórico e o progresso do rastreamento são perdidos, retornando o objeto ao ponto de origem.
    
    \item \textbf{Escalabilidade de Conexões:} O uso de SSE mantém uma conexão aberta por cliente. Em um cenário com um número muito elevado de usuários simultâneos, o gerenciamento dessas conexões pode exigir otimizações na infraestrutura do servidor.
    
    \item \textbf{Precisão Geográfica:} A simulação utiliza um modelo de movimento aleatório que não considera obstáculos físicos ou malhas viárias. O objeto pode se deslocar sobre áreas intransitáveis por não estar integrado a APIs de roteamento real.
    
    \item \textbf{Segurança e Autenticação:} Os canais de dados estão abertos para consulta pública, carecendo de camadas de autenticação ou controle de acesso para proteger as informações de localização.
\end{itemize}

\section{Evolução do Projeto}
Como etapa posterior, planeja-se transpor esta base de simulação para um cenário real: um sistema de rastreamento para motoboys de um restaurante local. A lógica de mensageria (SSE) e a estrutura de visualização em mapa serão reaproveitadas, substituindo a simulação matemática pela recepção de coordenadas GPS reais enviadas por dispositivos móveis dos entregadores.

\vfill
\begin{center}
    \textit{Este relatório foi estruturado para demonstrar o entendimento de concorrência em sistemas distribuídos e comunicação cliente-servidor de baixa latência.}
\end{center}

\end{document}
